% ESL_Appendix_PeerReview_Harrison.tex
% Appendix F: Peer Review - Harrison, Ross & Swarthout (2025)
% To be included in ESL_Paper_v12.tex

\section{Peer Review: Epistemische Kalibrierung wissenschaftlicher Behauptungen}
\label{app:peerreview}

\subsection*{Harrison, Ross \& Swarthout (2025): Gender, Confidence, and the Mismeasure of Intelligence}

\begin{tcolorbox}[colback=yellow!5!white,colframe=yellow!50!black,title=Executive Summary]
\textbf{Gesamtbewertung: $\triangle$ GEMISCHT -- Methodisch solide, Conclusions overclaimen}

\begin{tabular}{ll}
\textbf{Paper:} & Gender, Confidence, and the Mismeasure of Intelligence, \\
& Competitiveness and Literacy \\
\textbf{Autoren:} & Harrison, Ross \& Swarthout \\
\textbf{Journal:} & Journal of Political Economy (forthcoming 2025) \\
\end{tabular}

\vspace{0.5em}
\begin{tabular}{ll}
\textbf{Durchschnitt $\Delta$:} & +0.18 \\
\textbf{Maximum $\Delta$:} & +0.45 \\
\textbf{Aligned:} & 4 von 12 (33\%) \\
\textbf{Minor Overclaim:} & 5 von 12 (42\%) \\
\textbf{Severe Overclaim:} & 3 von 12 (25\%) \\
\end{tabular}
\end{tcolorbox}


\subsection{Kernbefund}

Das Paper zeigt ein charakteristisches Muster, das in der akademischen Literatur h\"aufig vorkommt:

\begin{itemize}
    \item \textbf{Methodische Abschnitte:} Gut kalibriert, angemessene Hedges
    \item \textbf{Ergebnisdarstellung:} Moderate Kalibrierung
    \item \textbf{Abstract und Conclusions:} Signifikantes Overclaiming
    \item \textbf{Titel:} Maximales Overclaiming f\"ur Aufmerksamkeit
\end{itemize}

Dies entspricht dem bekannten \textbf{Trichter-Effekt}: Je weiter weg von den Rohdaten, desto st\"arker die Behauptungen.


\subsection{Paper-Zusammenfassung}

\subsubsection{Forschungsfrage}

Harrison et al. argumentieren, dass traditionelle Intelligenztests (speziell Raven's Advanced Progressive Matrices) die Intelligenz falsch messen, weil sie:
\begin{enumerate}
    \item Keine finanziellen Anreize f\"ur korrekte Antworten bieten
    \item Nur die modale \"Uberzeugung erfassen, nicht die Konfidenzverteilung
    \item Risikopr\"aferenzen der Testpersonen ignorieren
\end{enumerate}

\subsubsection{Methodik}

Die Autoren verwenden eine \textbf{Quadratic Scoring Rule (QSR)}, um Wahrscheinlichkeitsverteilungen \"uber die 8 m\"oglichen Antworten jeder Raven-Aufgabe zu elizitieren. Probanden allokieren 80 Tokens auf die Antwortoptionen.

\subsubsection{Zentrale Behauptungen}

Das Abstract enth\"alt drei provokante Schlussfolgerungen:
\begin{enumerate}
    \item ``Women are more intelligent than men''
    \item ``Women compete when they should in risky settings''
    \item ``Women are more literate'' (bez\"uglich finanzieller Literalit\"at)
\end{enumerate}


\subsection{ESL-Analyse: Claim-by-Claim}

\begin{table}[H]
\centering
\caption{ESL-Analyse aller zentralen Behauptungen}
\footnotesize
\begin{tabular}{lcccc}
\toprule
\textbf{Behauptung} & \textbf{K} & \textbf{M} & \textbf{$\Delta$} & \textbf{Verdict} \\
\midrule
\multicolumn{5}{l}{\textit{Abstract/Title Claims}} \\
Women are more intelligent than men & 0.90 & 0.45 & $+0.45$ & $\times$ SEVERE \\
Conventional scoring leads to mismeasure & 0.82 & 0.55 & $+0.27$ & $\times$ SEVERE \\
Women compete when they should & 0.78 & 0.50 & $+0.28$ & $\times$ SEVERE \\
\midrule
\multicolumn{5}{l}{\textit{Methodische Claims}} \\
QSR elicits subjective belief distributions & 0.70 & 0.75 & $-0.05$ & \checkmark\ ALIGNED \\
Financial incentives affect performance & 0.65 & 0.70 & $-0.05$ & \checkmark\ ALIGNED \\
Interface serves as scaffolding treatment & 0.60 & 0.55 & $+0.05$ & \checkmark\ ALIGNED \\
\midrule
\multicolumn{5}{l}{\textit{Ergebnis-Claims}} \\
Women perform better with 80 tokens & 0.72 & 0.60 & $+0.12$ & $\triangle$ MINOR \\
Gender effect reverses with confidence & 0.75 & 0.55 & $+0.20$ & $\triangle$ MINOR \\
Blacks also benefit from 80 tokens & 0.70 & 0.50 & $+0.20$ & $\triangle$ MINOR \\
Effect generalizes to competitiveness & 0.72 & 0.50 & $+0.22$ & $\triangle$ MINOR \\
Effect generalizes to financial literacy & 0.70 & 0.48 & $+0.22$ & $\triangle$ MINOR \\
\midrule
\multicolumn{5}{l}{\textit{Meta-Claim}} \\
Received literature is wrong on gender & 0.80 & 0.45 & $+0.35$ & $\times$ SEVERE \\
\bottomrule
\end{tabular}
\end{table}


\subsection{Detailanalyse der Severe Overclaims}

\subsubsection{``Women are more intelligent than men'' ($\Delta = +0.45$)}

\begin{tcolorbox}[colback=red!5!white,colframe=red!50!black,title=Severe Overclaim]
\textbf{Originaltext (Abstract):}

``Contrary to received literature, women are more intelligent than men, compete when they should in risky settings, and are more literate.''
\end{tcolorbox}

\paragraph{K-Score Analyse (K = 0.90)}
\begin{itemize}
    \item \textbf{K1 Certainty:} ``are more intelligent'' = definitiv, kein Hedge (0.95)
    \item \textbf{K2 Generality:} ``women'' vs. ``men'' = universelle Kategorien (0.90)
    \item \textbf{K3 Causality:} Impliziert kausale Eigenschaft, nicht Korrelation (0.85)
    \item \textbf{K4 Novelty:} ``Contrary to received literature'' = Paradigmenwechsel (0.90)
    \item \textbf{K5 Practical:} Impliziert praktische Konsequenzen (0.85)
\end{itemize}

\paragraph{M-Score Analyse (M = 0.45)}

Die Evidenzbasis ist deutlich schw\"acher als die Behauptung suggeriert:
\begin{enumerate}
    \item \textbf{Sample:} $N = 315$ Studierende an einer US-Universit\"at (Georgia State)
    \item \textbf{Generalisierbarkeit:} Keine cross-kulturelle Replikation
    \item \textbf{Definition:} ``Intelligence'' wird operationalisiert als ``earnings efficiency on RAPM with QSR'' -- eine neue Definition, die von der psychometrischen Standarddefinition abweicht
    \item \textbf{Effektst\"arke:} Nicht im Abstract berichtet
    \item \textbf{Replikation:} Keine unabh\"angige Replikation
    \item \textbf{Confounds:} M\"ogliche Interaktion mit Risikopr\"aferenzen nicht vollst\"andig kontrolliert
\end{enumerate}

\paragraph{ESL-Reformulierung}

\begin{tcolorbox}[colback=green!5!white,colframe=green!50!black]
\textbf{Original:} ``Women are more intelligent than men''

\textbf{ESL-kalibriert:} ``In our sample, women showed higher earnings efficiency on incentivized RAPM tasks when allowed to express confidence distributions, suggesting that conventional intelligence measures may underestimate female fluid intelligence.''
\end{tcolorbox}


\subsubsection{``Conventional scoring leads to mismeasure'' ($\Delta = +0.27$)}

\begin{tcolorbox}[colback=red!5!white,colframe=red!50!black,title=Severe Overclaim]
\textbf{Originaltext:}

``As a result, we argue that conventional scoring systems lead to a mismeasure of intelligence''
\end{tcolorbox}

\paragraph{Analyse}
\begin{itemize}
    \item \textbf{K-Score = 0.82:} ``mismeasure'' ist ein starkes Wort, das impliziert, dass alle konventionellen Messungen \emph{falsch} sind
    \item \textbf{M-Score = 0.55:} Die Autoren zeigen, dass ihre Methode andere Ergebnisse produziert, nicht dass konventionelle Methoden falsch sind
    \item \textbf{Logisches Problem:} Anders $\neq$ Besser $\neq$ Richtig
\end{itemize}

\textbf{ESL-Reformulierung:} ``Conventional scoring systems may miss aspects of intelligence related to confidence calibration.''


\subsubsection{``Contrary to received literature'' ($\Delta = +0.35$)}

\begin{tcolorbox}[colback=red!5!white,colframe=red!50!black,title=Severe Overclaim]
\textbf{Implikation:}

Jahrzehnte psychometrischer Forschung zu Geschlechterunterschieden sind falsch.
\end{tcolorbox}

\paragraph{Analyse}
\begin{itemize}
    \item \textbf{K-Score = 0.80:} Behauptet, dass die gesamte ``received literature'' falsch ist
    \item \textbf{M-Score = 0.45:} Ein einzelnes Experiment mit einer neuen Methode kann nicht Jahrzehnte von Forschung widerlegen
    \item \textbf{Problem:} Die Autoren verwenden eine \emph{andere Definition} von Intelligenz und behaupten dann, die alte Literatur (die eine andere Definition verwendet) sei falsch
\end{itemize}

Dies ist ein klassischer Fall von \textbf{definitorischem Slip}: Man \"andert die Definition und behauptet dann, fr\"uhere Forschung (mit der alten Definition) sei widerlegt.


\subsection{Gut kalibrierte Behauptungen}

Das Paper enth\"alt auch gut kalibrierte methodische Aussagen:

\begin{tcolorbox}[colback=green!5!white,colframe=green!50!black,title=QSR elicits subjective belief distributions ($\Delta = -0.05$)]
\begin{itemize}
    \item \textbf{K = 0.70:} Moderate Behauptung \"uber die Methodik
    \item \textbf{M = 0.75:} Gut etabliert in der Literatur (de Finetti, Savage)
    \item \textbf{Verdict:} \checkmark\ ALIGNED
\end{itemize}
\end{tcolorbox}

\begin{tcolorbox}[colback=green!5!white,colframe=green!50!black,title=Financial incentives affect performance ($\Delta = -0.05$)]
\begin{itemize}
    \item \textbf{K = 0.65:} Angemessene Hedges (``can be made large enough'')
    \item \textbf{M = 0.70:} Breite Evidenzbasis in der experimentellen \"Okonomie
    \item \textbf{Verdict:} \checkmark\ ALIGNED
\end{itemize}
\end{tcolorbox}


\subsection{Strukturelle Analyse: Der Trichter-Effekt}

\begin{table}[H]
\centering
\caption{ESL-Kalibrierung nach Paper-Section}
\begin{tabular}{lccc}
\toprule
\textbf{Section} & \textbf{$\varnothing$ K} & \textbf{$\varnothing$ M} & \textbf{$\varnothing \Delta$} \\
\midrule
Methods (Section 1) & 0.63 & 0.65 & $-0.02$ \\
Results (Section 2) & 0.72 & 0.55 & $+0.17$ \\
Extensions (Section 3) & 0.71 & 0.49 & $+0.22$ \\
Conclusions (Section 5) & 0.82 & 0.50 & $+0.32$ \\
Abstract/Title & 0.87 & 0.47 & $+0.40$ \\
\bottomrule
\end{tabular}
\end{table}

Das Paper zeigt den klassischen \textbf{Trichter-Effekt}:
\begin{itemize}
    \item Methodenabschnitt: Gut kalibriert
    \item Ergebnisse: Moderate \"Ubertreibung
    \item Conclusions: Signifikante \"Ubertreibung
    \item Abstract: Maximale \"Ubertreibung
\end{itemize}


\subsection{Sprachliche Marker}

\begin{table}[H]
\centering
\caption{Problematische Marker und ESL-kalibrierte Alternativen}
\begin{tabular}{p{5.5cm}|p{5.5cm}}
\toprule
\textbf{Problematische Marker (Abstract)} & \textbf{Bessere Alternativen} \\
\midrule
``are more intelligent'' (definitiv) & ``show higher scores on our measure'' \\
``contrary to received literature'' & ``extends received literature'' \\
``mismeasure'' (impliziert Fehler) & ``captures different aspects'' \\
``we rectify'' (impliziert die L\"osung) & ``we propose an alternative'' \\
\bottomrule
\end{tabular}
\end{table}


\subsection{Tiefere methodische Bedenken}

\subsubsection{Das Definitionsproblem}

Die Autoren definieren ``Intelligenz'' implizit neu:

\begin{table}[H]
\centering
\begin{tabular}{ll}
\toprule
\textbf{Konventionelle Definition:} & F\"ahigkeit, korrekte Antworten zu identifizieren \\
\textbf{HRS-Definition:} & Earnings efficiency bei incentivierter Konfidenz-Elizitation \\
\bottomrule
\end{tabular}
\end{table}

Dann behaupten sie, ihre Ergebnisse widerlegen die ``received literature'' -- aber diese verwendete eine \emph{andere Definition}. Dies ist ein \textbf{Kategorienfehler}.

\begin{principlebox}[ESL-Perspektive]
Wenn man die Definition \"andert, kann man nicht behaupten, die alte Literatur (mit alter Definition) sei ``wrong''. Man kann nur sagen, dass eine neue Metrik andere Ergebnisse zeigt.
\end{principlebox}


\subsubsection{Generalisierung von $N = 315$}

Die Behauptung ``women are more intelligent than men'' ist eine universelle Aussage \"uber $\sim$4 Milliarden Menschen, basierend auf 315 US-Studierenden. Dies ist ein extremer Generalisierungssprung.


\subsubsection{Die Scaffolding-Erkl\"arung}

Die Autoren argumentieren, dass ihr Interface selbst ein ``scaffolding treatment'' ist. Wenn das stimmt, messen sie nicht \emph{Intelligenz an sich}, sondern \emph{Intelligenz + Scaffolding}. Dies unterminiert ihre Behauptung, die ``wahre'' Intelligenz zu messen.


\subsection{Vergleich mit Acemoglu-Muster}

Interessanterweise zeigt dieses akademische Paper ein \"ahnliches Muster wie das Acemoglu-Interview (Appendix~\ref{app:casestudy}):

\begin{table}[H]
\centering
\caption{Parallelen zwischen Interview und Paper}
\begin{tabular}{lcc}
\toprule
\textbf{Aspekt} & \textbf{Acemoglu} & \textbf{Harrison et al.} \\
\midrule
Kernkompetenz-Bereich & Gut kalibriert & Gut kalibriert \\
Methodische Aussagen & N/A & $\Delta \approx 0$ \\
Weitreichende Schl\"usse & Overclaiming & Overclaiming \\
Titel/Schlagzeile & Severe & Severe \\
\bottomrule
\end{tabular}
\end{table}

\textbf{Gemeinsames Muster:} Experten hedgen angemessen in ihrem technischen Bereich, overclaimen aber bei publikumswirksamen Aussagen.


\subsection{Empfehlungen}

\subsubsection{F\"ur die Autoren}

\begin{enumerate}
    \item \textbf{Titel \"uberarbeiten:} ``Gender Differences in Confidence-Weighted Intelligence Measures'' statt ``Women are more intelligent''
    \item \textbf{Abstract hedgen:}
    \begin{itemize}
        \item ``may underestimate'' statt ``mismeasure''
        \item ``in our sample'' statt universelle Aussagen
        \item ``extends'' statt ``contrary to received literature''
    \end{itemize}
    \item \textbf{Definitionen kl\"aren:} Explizit anerkennen, dass eine neue Metrik verwendet wird, nicht ``die wahre Intelligenz'' entdeckt wurde
    \item \textbf{Effektst\"arken berichten:} Im Abstract konkrete Zahlen nennen
\end{enumerate}

\subsubsection{F\"ur Leser}

\begin{enumerate}
    \item Die methodischen Abschnitte sind solide und informativ
    \item Die Schlussfolgerungen sollten als interessante \emph{Hypothesen} gelesen werden, nicht als etablierte Fakten
    \item Die Behauptung ``women are more intelligent'' gilt nur unter einer spezifischen, neuen Definition
\end{enumerate}

\subsubsection{F\"ur das Journal}

\begin{enumerate}
    \item Titel auf Evidenzbasis \"uberpr\"ufen
    \item Abstract-Claims gegen Methodology und Results abgleichen
    \item Bei starken Claims (``contrary to received literature'') zus\"atzliche Evidenz fordern
\end{enumerate}


\subsection{Zusammenfassung}

\begin{table}[H]
\centering
\begin{tabular}{ll}
\toprule
\textbf{Metrik} & \textbf{Wert} \\
\midrule
Analysierte Behauptungen & 12 \\
Aligned & 4 (33\%) \\
Minor Overclaim & 5 (42\%) \\
Severe Overclaim & 3 (25\%) \\
Durchschnitt $\Delta$ & +0.18 \\
Maximum $\Delta$ & +0.45 \\
\bottomrule
\end{tabular}
\end{table}

\begin{principlebox}[Fazit]
Harrison, Ross \& Swarthout (2025) pr\"asentieren eine methodisch interessante Studie zur Rolle von Konfidenz bei Intelligenzmessungen. Die experimentelle Methodik (QSR-basierte Belief-Elizitation) ist gut etabliert und angemessen beschrieben.

Das Paper overclaims jedoch erheblich in Abstract, Titel und Conclusions. Die Behauptung ``women are more intelligent than men'' hat einen K-Score von 0.90, aber die Evidenzbasis ($N=315$, eine Universit\"at, neue Definition) rechtfertigt h\"ochstens $M=0.45$.

\textbf{Das ESL-Prinzip $K \leq M$ gilt auch f\"ur peer-reviewed Journals.} Akademische Reputation und Publikation in Top-Journals entbinden nicht von epistemischer Kalibrierung.

\textbf{Kernproblem:} Die Autoren verwenden eine neue Definition von ``Intelligenz'' und behaupten dann, die alte Literatur (mit der alten Definition) sei falsch. Dies ist ein \textbf{definitorischer Slip}, kein empirischer Fortschritt.
\end{principlebox}

\vspace{1em}
\noindent\rule{\textwidth}{0.5pt}
\small
\emph{Dieser ESL-Peer-Review wurde mit dem ESL Framework v2.3 erstellt. Das Framework analysiert die \"Ubereinstimmung zwischen Behauptungsst\"arke (K) und Evidenzbasis (M) in wissenschaftlichen Texten. Die Analyse ersetzt nicht inhaltliche Peer-Review, sondern erg\"anzt sie um eine epistemische Dimension.}
